\documentclass[11pt,reqno]{amsart}
\usepackage[localtheorems]{raeez-math-template}

% ================================================================
% Packages
% ================================================================
\usepackage{array}

% ================================================================
% Theorem environments
% ================================================================
\newtheorem{theorem}{Theorem}[section]
\newtheorem{proposition}[theorem]{Proposition}
\newtheorem{lemma}[theorem]{Lemma}
\newtheorem{corollary}[theorem]{Corollary}
\theoremstyle{definition}
\newtheorem{definition}[theorem]{Definition}
\newtheorem{example}[theorem]{Example}
\theoremstyle{remark}
\newtheorem{remark}[theorem]{Remark}

% ================================================================
% Macros
% ================================================================
\newcommand{\cA}{\mathcal{A}}
\newcommand{\cB}{\mathcal{B}}
\newcommand{\cC}{\mathcal{C}}
\newcommand{\cD}{\mathcal{D}}
\newcommand{\cW}{\mathcal{W}}
\newcommand{\barB}{\overline{B}}
\newcommand{\barC}{\overline{C}}
\newcommand{\Ran}{\mathrm{Ran}}
\newcommand{\MC}{\mathrm{MC}}
\newcommand{\Sym}{\mathrm{Sym}}
\newcommand{\Hom}{\mathrm{Hom}}
\newcommand{\Ext}{\mathrm{Ext}}
\newcommand{\Res}{\mathrm{Res}}
\newcommand{\Aut}{\mathrm{Aut}}
\newcommand{\fg}{\mathfrak{g}}
\newcommand{\Ainf}{A_\infty}
\newcommand{\Vir}{\mathrm{Vir}}
\newcommand{\Conv}{\mathrm{Conv}}

\DeclareMathOperator{\gr}{gr}
\DeclareMathOperator{\rk}{rk}
\DeclareMathOperator{\wt}{wt}

\numberwithin{equation}{section}

% ================================================================
\begin{document}

\title[Completed chiral bar-cobar duality and MC4 weight cutoff]%
{Completed chiral bar-cobar duality and MC4 weight cutoff}

\author{Raeez Lorgat}
\address{Department of Mathematics, Harvard University,
 Cambridge, MA 02138, USA}
\email{rlorgat@math.harvard.edu}

\date{}

\begin{abstract}
The classical bar-cobar adjunction of Priddy produces a Koszul
dual for a quadratic algebra. For non-quadratic chiral algebras
on a curve, completed bar constructions require two independent
inputs: a degree cutoff for the Maurer-Cartan equation and a
Mittag-Leffler condition for the inverse system of finite-stage
cohomologies. A \emph{strong completion tower} is a filtered
augmented curved chiral $\Ainf$-algebra whose operations are
filtration-nondecreasing on the augmentation ideal. Under this
axiom the Maurer-Cartan equation is componentwise finite. If the
finite-stage cohomology tower is Mittag-Leffler, the completed
bar-cobar comparison is a quasi-isomorphism in the square-zero
total model; in the genuinely curved CDG ambient it is a
coacyclic-cone comparison in the second-kind coderived category.

The standard conformal filtrations on affine Kac-Moody, Virasoro,
principal $\cW$-algebras, and lattice vertex algebras do not by
themselves supply the algebra-level strong filtration axiom.
They enter through bar-level weight-cutoff criteria. The
positive MC4${}^+$ regime is the $\rho=0$ case, proved by strong
completion or by a genuine auxiliary positive grading. The generic
finite-resonance MC4${}^0$ regime, including Virasoro and
principal $\cW_N$ at generic parameters, is proved on the
positive-energy weight-completed finite-resonance surface. A bare
screening presentation still requires the explicit chain-level
strong deformation retract or an independent stabilization theorem
that identifies it with that surface.
\end{abstract}

\subjclass[2020]{17B69, 18M70, 14H10, 81T40}

\keywords{Chiral algebra, bar construction, Koszul duality,
Maurer-Cartan, curved $A_\infty$-algebra, completion,
Mittag-Leffler, $W$-algebra}

\maketitle

% ================================================================
% 1. INTRODUCTION
% ================================================================

\section{Introduction}\label{sec:intro}

\subsection{The quadratic case and its limits}

Classical Koszul duality is a phenomenon of graded algebras.
A quadratic algebra $A = T(V)/(R)$ determines the dual
coalgebra $A^{\mathrm i}\hookrightarrow T^c(sV^*)$, the maximal
subcoalgebra whose quadratic core is $s^2R^\perp$; the bar
construction $B(A)$ mediates between them; and $A$ is Koszul when
$H^*(B(A))$ concentrates in bar length one, so that
$\Omega(A^{\mathrm i}) \simeq A$. The theory is due to
Priddy~\cite{Priddy70}, Beilinson--Ginzburg--Soergel~\cite{BGS},
and in
operadic form Loday-Vallette~\cite{LV12}.

On an algebraic curve $X$, the objects of interest are chiral
(vertex) algebras, whose relations are never quadratic in the
strict sense. Heisenberg has a central term in the
$J \cdot J$ OPE. Virasoro has a quartic pole. Principal
$\cW$-algebras $\cW^k(\fg, f_{\mathrm{prin}})$ have generators
of weights $m_i+1$, where $m_i$ are the exponents of~$\fg$;
only for $\fg=\mathfrak{sl}_N$ are the weights $2,\ldots,N$.
Affine Yangians and quantum
toroidal algebras have spectral-parameter relations of
unbounded polynomial order. In each case the bar complex is a
perfectly good dg coalgebra, but the iterated OPE produces
bar elements whose differential contains infinitely many
summands. A finite-dimensional truncation sees only a finite
part of this object, and the classical bar-cobar adjunction is
silent.

\subsection{The completion problem}

The remedy is completion. One introduces a decreasing
filtration
$\cA = F^0\cA \supset F^1\cA \supset \cdots$
with finite-dimensional quotients
$\cA_{\le N} := \cA/F^{N+1}\cA$ and writes
$\widehat{\barB}(\cA) := \varprojlim_N \barB(\cA_{\le N})$.
At each finite stage $N$ the strict bar-cobar adjunction
produces a universal twisting morphism
$\tau_N \colon \barB(\cA_{\le N}) \to \cA_{\le N}$
which is a Maurer-Cartan element of the convolution dg Lie
algebra
$\Conv(\barB(\cA_{\le N}), \cA_{\le N})$. The completion
problem, labelled MC4 in~\cite{LorgatI}, asks: do the
finite-stage MC elements $\tau_N$ glue to a single MC
element $\widehat\tau$ in the completed convolution
algebra, and is the twisted tensor product acyclic in the
limit?

Two obstructions appear. First, the MC equation
$\partial(\widehat\tau) + \widehat\tau \star \widehat\tau = 0$
is a sum over all bar degrees, and a priori each component is
infinite. Second, the passage from finite-stage acyclicity to
acyclicity in the limit requires the Mittag-Leffler condition,
which is a separate finite-window hypothesis.

\subsection{The main result}

The central structural axiom is that the $\Ainf$ operations of
$\cA$ respect the additive weight:
\begin{equation}\label{eq:SF-intro}
\mu_r(F^{i_1}\cA, \dots, F^{i_r}\cA) \subset F^{i_1 + \cdots + i_r}\cA.
\end{equation}
A filtered curved chiral $\Ainf$-algebra for which
\eqref{eq:SF-intro} holds and whose augmentation ideal
$\overline{\cA} = F^1\cA$ is the first filtration piece is a
\emph{strong completion tower}
(Definition~\ref{def:strong-tower}).

\begin{theorem}[Strong completion theorem, informal form]
Let $\cA$ be a strong completion tower whose finite-stage
quotients lie in the proved square-zero bar-cobar regime and whose
finite-stage cohomology tower is Mittag-Leffler. Then the
completed bar construction $\widehat{\barB}(\cA)$ is a
separated complete pronilpotent curved dg chiral coalgebra,
the completed convolution algebra carries a canonical
Maurer-Cartan element $\widehat\tau = \varprojlim_N \tau_N$,
and the counit
\[
\widehat{\epsilon}\colon
\widehat{\Omega}\bigl(\widehat{\barB}(\cA)\bigr)
\xrightarrow{\ \sim\ }\cA
\]
is a quasi-isomorphism in the square-zero total model. In the
curved CDG ambient the same finite-stage construction gives a
coacyclic-cone comparison under the second-kind hypotheses.
\end{theorem}

The precise statement is Theorem~\ref{thm:mc4-strong} below.
The proof uses three ingredients:
(i) the degree cutoff lemma (\S\ref{sec:cutoff}), which reduces
each component of the MC equation to a finite sum;
(ii) the finite-window Mittag-Leffler criterion
(\S\ref{sec:ml}); and (iii) the reduction principle of
Lemma~\ref{lem:filtered-comparison} which turns the finite-stage
theory into a limit statement via the Milnor exact sequence.

\subsection{Positive versus resonant towers}

Completion problems split into positive and resonant regimes. A
tower is \emph{positive} if the
filtration is defined by a genuine positive grading (e.g.\ the
$L_0$-eigenvalue filtration on a graded vertex algebra). The
degree cutoff lemma then holds after the strong axiom or the
bar-level weight-cutoff criterion is verified. This is the
MC4${}^+$ regime.

A tower is \emph{resonant} if there is a finite-dimensional
weight-zero piece $R_\cA$ on which higher operations can
preserve filtration degree. The \emph{resonance rank}
\[
\rho(\cA) := \dim H^*(R_\cA, m_1^{R_\cA})
\]
classifies the completion difficulty. Positive towers have
$\rho = 0$; Virasoro and the principal $\cW_N$-algebras have
$0 < \rho < \infty$; genuinely wild algebras have
$\rho = \infty$. The MC4${}^0$ problem concerns the middle
regime. For Virasoro and principal $\cW_N$ at generic parameters,
the positive-energy weight-completed finite-resonance model proves
the completed comparison. A concrete screening presentation of a
principal $\cW_N$ algebra enters this theorem only after one gives
the chain-level strong deformation retract, or an independent
weight-cutoff stabilization theorem, identifying it with that
finite-resonance model.

\subsection{Strong filtration in curved bar-cobar theory}

Positselski~\cite{Positselski11} developed curved bar
constructions for associative algebras with curvature.
Loday-Vallette~\cite{LV12} supply the operadic bar-cobar
machinery on which the chiral lift is built. The
Beilinson-Drinfeld framework of filtered chiral algebras with
$I$-adic topology~\cite{BD04} provides the ambient topology
for the inverse limit. The theorem below is the chiral and
$\Ainf$ lift: the novelty is the strong filtration axiom,
which identifies the precise structural hypothesis under
which all finite-stage data assemble without further input.

% ================================================================
% 2. PRELIMINARIES
% ================================================================

\section{The bar construction on the augmentation ideal}
\label{sec:prelim}

Throughout, $X$ is a smooth complex algebraic curve and
$\cA$ is an augmented curved chiral $\Ainf$-algebra on $X$ in
the sense of the monograph~\cite{LorgatI}. Augmentation means
that $\cA$ is equipped with a counit
$\varepsilon\colon \cA \to \mathbf{1}$ splitting the unit.
The augmentation ideal is
\[
\overline{\cA} := \ker(\varepsilon).
\]

\begin{remark}[Augmentation ideal]
\label{rem:barA}
The bar construction is built on $\overline{\cA}$, not on
$\cA$: writing $T^c(s^{-1}\cA)$ in place of
$T^c(s^{-1}\overline{\cA})$ includes the unit and breaks the
Koszul sign computation. We therefore set
\[
\barB^{\mathrm{ch}}(\cA) := T^c\bigl(s^{-1}\overline{\cA}\bigr),
\]
with $s^{-1}$ the homological desuspension
$|s^{-1}v| = |v| - 1$ and with the chiral coproduct obtained
by deconcatenation on the collision divisor of
$\overline{C}_{\bullet+1}(X)$.
\end{remark}

The curved chiral $\Ainf$ operations $\{\mu_r\}_{r \ge 0}$
assemble into a codifferential $d_{\barB}$ on
$\barB^{\mathrm{ch}}(\cA)$ with curvature coming from $\mu_0$,
so $\barB^{\mathrm{ch}}(\cA)$ is a curved dg coalgebra. The
convolution dg Lie algebra
$\Conv(\cC, \cA) := \Hom(\cC, \cA)$ has bracket from the
coproduct on $\cC$ and the product on $\cA$; a Maurer-Cartan
element is a degree-$1$ map $\tau\colon \cC \to \cA$ with
$\partial(\tau) + \tau \star \tau = 0$.

\begin{theorem}[Finite-type bar-cobar duality, cited]
\label{thm:finite-type-bar-cobar}
Let $\cA$ be a finite-type augmented curved chiral
$\Ainf$-algebra on $X$ in the proved square-zero bar-cobar regime.
Then the universal twisting morphism
$\tau\colon \barB^{\mathrm{ch}}(\cA) \to \cA$ is a
Maurer-Cartan element of
$\Conv(\barB^{\mathrm{ch}}(\cA), \cA)$, and the counit
\[
\epsilon\colon \Omega^{\mathrm{ch}}(\barB^{\mathrm{ch}}(\cA))
\xrightarrow{\ \sim\ }\cA
\]
is a quasi-isomorphism in the square-zero total model. In the
curved CDG ambient the corresponding finite-stage cone is
coacyclic in the second-kind coderived category.
\end{theorem}

\begin{proof}
See~\cite[Thm.~B]{LorgatI}. The finite-type hypothesis
reduces to finite-type vertex algebras, where the operadic
bar-cobar machinery of Loday-Vallette~\cite{LV12} applies
after the curved extension of Positselski~\cite{Positselski11}.
\end{proof}

% ================================================================
% 3. THE STRONG FILTRATION AXIOM
% ================================================================

\section{The strong filtration axiom}
\label{sec:sf}

\begin{definition}[Strong completion tower]
\label{def:strong-tower}
An augmented curved chiral $\Ainf$-algebra $\cA$ on a curve
$X$ is a \emph{strong completion tower} if it carries a
descending filtration
\[
\cA = F^0\cA \supset F^1\cA \supset F^2\cA \supset \cdots,
\qquad
\bigcap_{N \ge 0} F^{N+1}\cA = 0,
\]
such that:
\begin{enumerate}
\item $\cA$ is separated and complete:
 $\cA \cong \varprojlim_N \cA_{\le N}$, where
 $\cA_{\le N} := \cA/F^{N+1}\cA$;
\item every quotient $\cA_{\le N}$ is of finite type and lies
 in the proved bar-cobar regime of
 Theorem~\ref{thm:finite-type-bar-cobar};
\item $\overline{\cA} = F^1\cA$: the augmentation ideal is
 the first filtration piece;
\item the curvature term satisfies $\mu_0(1)\in F^1\cA$, and all
	 positive-arity chiral $\Ainf$-operations are
	 filtration-nondecreasing,
	 \begin{equation}\label{eq:strong-filtration}
	 \mu_r(F^{i_1}\cA, \dots, F^{i_r}\cA) \subset
	 F^{i_1 + \cdots + i_r}\cA
	 \end{equation}
	 for $r\ge 1$.
\end{enumerate}
\end{definition}

Axiom~(3) ties the filtration to the augmentation. Axiom~(4)
is the structural heart of the definition: the curvature starts in
the augmentation ideal, and the positive-arity operations respect
the additive grading.

\begin{remark}[On the quotient maps]
Each projection $p_N\colon \cA_{\le N+1} \twoheadrightarrow
\cA_{\le N}$ is a strict morphism of curved chiral
$\Ainf$-algebras, because the kernel
$F^{N+1}\cA/F^{N+2}\cA$ is an $\Ainf$-ideal of
$\cA_{\le N+1}$ by~\eqref{eq:strong-filtration}. The
finite-stage bar constructions $\barB^{\mathrm{ch}}(\cA_{\le N})$
therefore form an inverse system of curved dg chiral
coalgebras with strict morphisms, and the completed bar
construction is the literal inverse limit.
\end{remark}

\begin{proposition}[Standard landscape weight cutoff]
\label{prop:examples}
For the standard conformal-weight filtration
$F^N\cA := \bigoplus_{h \ge N}\cA_h$, the following families
satisfy the bar-level weight-cutoff criterion: in every fixed
total conformal weight only finitely many generators and OPE
residue terms contribute to the completed bar differential.
\begin{enumerate}[label=\textup{(\alph*)}]
\item the affine Kac-Moody vertex algebra $V_k(\fg)$ for
 any simple $\fg$ and any non-critical level
 $k \ne -h^\vee$;
\item the Virasoro vertex algebra $\Vir_c$ for any $c$;
\item the principal $\cW$-algebra $\cW^k(\fg, f_{\mathrm{prin}})$
 for any simple $\fg$ and any non-critical $k$;
\item the lattice vertex algebra $V_\Lambda$ for any
 positive-definite even lattice $\Lambda$.
\end{enumerate}
This statement does not assert the algebra-level strong
filtration axiom
$\mu_r(F^{i_1},\ldots,F^{i_r})\subset F^{i_1+\cdots+i_r}$.
\end{proposition}

\begin{proof}
Each family is $L_0$-graded with finite-dimensional homogeneous
pieces and finitely many strong generators in each bounded weight
window. The OPE of two homogeneous fields of weights $a,b$ has
only finitely many poles, and every residue term has conformal
weight at most $a+b-1$. Thus a bar chain of fixed total weight
uses only fields of bounded weight and only finitely many residue
terms. This is the bar-level cutoff needed for the finite-window
MC4 comparison.

The assertion is weaker than algebra-level strong filtration.
For affine currents, for example, the term
$[J^a,J^b]$ has weight~$1$ although the inputs have weights
$1$ and~$1$; hence the strong axiom fails for the raw conformal
filtration. Principal $\cW^k(\fg,f_{\mathrm{prin}})$ is governed
by the exponents of~$\fg$: the generator weights are $m_i+1$,
with $m_i$ running through the exponents.
\end{proof}

\begin{example}[Non-example: $\beta\gamma$-system with
conformal grading]
\label{ex:non-example}
The $\beta\gamma$-system has a weight-$\tfrac12$ grading that
refines the conformal grading. With respect to the
conformal-weight filtration alone, the strong filtration
axiom fails because the OPE
$\beta(z)\gamma(w) \sim 1/(z-w)$ produces a weight-zero
output from weight-$\tfrac12$ inputs, violating additivity.
Its shadow tower is contact class~$C$, with a nonzero quartic
contact shadow; it is not a Heisenberg-type degree-$2$ tower.
A different auxiliary filtration may still give a bar-level
weight cutoff.
\end{example}

% ================================================================
% 4. THE DEGREE CUTOFF LEMMA
% ================================================================

\section{The degree cutoff lemma}
\label{sec:cutoff}

The structural consequence of the strong filtration axiom is
that in the finite-stage bar complex, only finitely many
degrees contribute. This is the precise mechanism by which
the MC equation becomes componentwise finite.

\begin{lemma}[Degree cutoff]
\label{lem:degree-cutoff}
Let $\cA$ be a strong completion tower and let
$\cA_{\le N} = \cA/F^{N+1}\cA$ be the $N$-th quotient. In the
finite-stage convolution algebra
$\Conv(\barB^{\mathrm{ch}}(\cA_{\le N}), \cA_{\le N})$, the
Maurer-Cartan equation
\[
\partial(\tau_N) + \tau_N \star \tau_N = 0
\]
involves only degrees $r \le N$: on $\cA_{\le N}$ the bar
differential is the finite sum $b_0+b_1+\cdots+b_N$
of coderivations coming from $\mu_0,\mu_1,\dots,\mu_N$.
\end{lemma}

\begin{proof}
The arity-zero curvature term $b_0$ is a single continuous term
because $\mu_0(1)\in F^1\cA$. Every positive-arity bar input lies
in $\overline{\cA}_{\le N}$, which by
axiom~(3) equals the image of $F^1\cA$ in the quotient. The
strong filtration axiom~\eqref{eq:strong-filtration} gives
\[
\mu_r(\overline{\cA}^{\otimes r}) \subset
\mu_r\bigl((F^1\cA)^{\otimes r}\bigr) \subset F^r\cA.
\]
Reducing modulo $F^{N+1}\cA$, any term with $r \ge N+1$ lands
in $F^r\cA \subset F^{N+1}\cA$ and vanishes in $\cA_{\le N}$.
So the induced bar differential on $\cA_{\le N}$ is the
finite sum
$b_0+b_1+\cdots+b_N$, as claimed.
\end{proof}

\begin{remark}[Degree cutoff as PBW degeneration]
The degree cutoff is the operadic counterpart of the classical
PBW theorem. For the universal enveloping algebra $U(\fg)$
with its standard filtration by number of generators, the
associated graded is $\mathrm{Sym}(\fg)$; higher-degree
brackets vanish modulo lower filtration stages. The strong
filtration axiom extends this phenomenon to all chiral
$\Ainf$-operations on a graded vertex algebra: the degree
cutoff lemma is the statement that at each finite stage,
only finitely many Poisson (or vertex) operations are
active. The classical PBW theorem is the $r = 2$ case.
\end{remark}

\begin{corollary}[Finite-stage MC assembly]
\label{cor:finite-stage-mc}
Let $\cA$ be a strong completion tower. The universal
twisting morphism $\tau_N$ at stage $N$ is a Maurer-Cartan
element of the finite-stage, locally finite convolution
$\Ainf$-algebra: the quotient
$\barB^{\mathrm{ch}}(\cA_{\le N})$ has finite-dimensional
bar-length components, and the bar differential is the finite
sum of Lemma~\ref{lem:degree-cutoff}.
\end{corollary}

\begin{proof}
Finite type of $\cA_{\le N}$ gives finite dimension of each
$\overline{\cA}_{\le N}^{\otimes r}$ for every $r$, and the
degree cutoff lemma bounds the number of nonzero bar
differential terms. Theorem~\ref{thm:finite-type-bar-cobar}
then gives the Maurer-Cartan property.
\end{proof}

% ================================================================
% 5. MITTAG-LEFFLER FINITE-WINDOW CRITERION
% ================================================================

\section{Mittag-Leffler Finite-Window Criterion}
\label{sec:ml}

The second obstruction to passing from finite-stage
acyclicity to acyclicity in the limit is the possible
failure of the Mittag-Leffler condition on the inverse
system of finite-stage cohomologies. The strong filtration
axiom supplies the degree cutoff; the Mittag-Leffler condition
is a separate finite-window hypothesis.

\begin{proposition}[Mittag-Leffler criterion for finite windows]
\label{prop:ml}
Let $\cA$ be a filtered tower whose finite-stage bar complexes
have finite weight windows: for each total weight $w$ there is
$N(w)$ such that the weight-$w$ subcomplex of
$\barB^{\mathrm{ch}}(\cA_{\le N})$ is independent of~$N$ for
$N\ge N(w)$. Then, for every cohomological degree $m$, the
inverse system
\[
\bigl\{H^m(\barB^{\mathrm{ch}}(\cA_{\le N}))\bigr\}_{N \ge 0}
\]
is Mittag-Leffler on every fixed weight slice. Consequently the
finite weight truncations, and the product completion over weight
slices, have vanishing $\varprojlim^1$.
\end{proposition}

\begin{proof}
The transition maps are surjections because the finite-stage bar
construction is functorial on surjections of finite-type curved
chiral $\Ainf$-algebras. By hypothesis, for each fixed weight
$w$ there is $N(w)$ such that the transition map
\[
\barB^{\mathrm{ch}}(\cA_{\le N+1})_w \xrightarrow{\ \sim\ }
\barB^{\mathrm{ch}}(\cA_{\le N})_w
\]
is an isomorphism of subcomplexes for $N\ge N(w)$. Passing to
cohomology gives stable isomorphisms on the weight-$w$ slice.

For a finite weight truncation only finitely many slices occur,
so the inverse system is Mittag-Leffler. For the completed
product over all weights, inverse limits and products are taken
slicewise over vector spaces; the preceding stabilization kills
the derived limit on each slice, hence on the completed product.
\end{proof}

\begin{lemma}[Complete filtered comparison]
\label{lem:filtered-comparison}
Let $\{f_N\colon C_N \to D_N\}$ be a morphism of inverse
systems of complexes with surjective transition maps. If each
$f_N$ is a quasi-isomorphism and the inverse systems
$\{H^*(C_N)\}$ and $\{H^*(D_N)\}$ are Mittag-Leffler, then
the induced map on limits
$\widehat f\colon \widehat C \to \widehat D$
is a quasi-isomorphism.
\end{lemma}

\begin{proof}
For any inverse system of cochain complexes the Milnor
exact sequence reads
\[
0 \to \varprojlim\nolimits^1_N H^{m-1}(C_N) \to
H^m(\widehat C) \to \varprojlim_N H^m(C_N) \to 0.
\]
Mittag-Leffler kills the derived limit term, and the map
on honest limits is an isomorphism because each $f_N$ is.
\end{proof}

\section{The main theorem}
\label{sec:main}

\begin{theorem}[Strong completion tower bar-cobar duality]
\label{thm:mc4-strong}
Let $\cA$ be a strong completion tower
(Definition~\ref{def:strong-tower}) whose finite-stage comparison
systems satisfy the finite-window Mittag-Leffler condition of
Proposition~\ref{prop:ml}. Write
$\tau_N \in \MC\bigl(\Conv(\barB^{\mathrm{ch}}(\cA_{\le N}),
\cA_{\le N})\bigr)$ for the universal twisting morphism at
stage $N$. Then:
\begin{enumerate}
\item \emph{Completed coalgebra.} The completed bar
 construction
 \[
 \widehat{\barB}^{\mathrm{ch}}(\cA) := \varprojlim_N
 \barB^{\mathrm{ch}}(\cA_{\le N})
 \]
 exists as a separated complete pronilpotent curved dg
 chiral coalgebra with continuous differential, coproduct,
 and curvature. At each stage,
 $\widehat{\barB}^{\mathrm{ch}}(\cA)/F^{N+1} \cong
 \barB^{\mathrm{ch}}(\cA_{\le N})$.
\item \emph{Completed MC element.} The componentwise limit
 $\widehat\tau = \varprojlim_N \tau_N$ is a Maurer-Cartan
 element of the completed convolution algebra
 \[
 \widehat\Conv := \varprojlim_N
 \Conv(\barB^{\mathrm{ch}}(\cA_{\le N}), \cA_{\le N}),
 \]
 and the MC equation
 $\partial(\widehat\tau) + \widehat\tau \star \widehat\tau
 = 0$
 converges in the inverse-limit topology, each component
 being a finite sum by the degree cutoff lemma.
\item \emph{Completed cobar.} The completed cobar object
 \[
 \widehat{\Omega}^{\mathrm{ch}}\bigl(
 \widehat{\barB}^{\mathrm{ch}}(\cA)\bigr) :=
 \varprojlim_N
 \Omega^{\mathrm{ch}}(\barB^{\mathrm{ch}}(\cA_{\le N}))
 \]
 is a separated complete curved dg chiral algebra, equal to
 the twisted tensor product
 $\cA \otimes_{\widehat\tau} \widehat{\barB}^{\mathrm{ch}}(\cA)$.
\item \emph{Counit in the square-zero total model.} The completed
	 counit
 \[
 \widehat\epsilon \colon
 \widehat{\Omega}^{\mathrm{ch}}\bigl(
 \widehat{\barB}^{\mathrm{ch}}(\cA)\bigr)
 \xrightarrow{\ \sim\ }\cA
 \]
	 is a quasi-isomorphism when both sides carry the
	 square-zero total differential. In the genuinely curved CDG
	 ambient, the same comparison is replaced by a coacyclic-cone
	 statement in the specified second-kind coderived category.
\item \emph{Dual comparison.} If
	 $\cC = \varprojlim_N \cC_{\le N}$ is a separated complete
	 pronilpotent curved dg chiral coalgebra whose finite
	 quotients lie in the finite-type square-zero bar-cobar
	 regime and whose finite-window tower is Mittag-Leffler, then the
	 completed unit
 \[
 \widehat\eta\colon \cC \xrightarrow{\ \sim\ }
 \widehat{\barB}^{\mathrm{ch}}\bigl(
 \widehat{\Omega}^{\mathrm{ch}}(\cC)\bigr)
 \]
	 is a quasi-isomorphism in the square-zero total model, with
	 the curved CDG version read as a coacyclic-cone comparison.
\item \emph{Equivalence surface.} Completed bar and completed
	 cobar give quasi-inverse equivalences on the homotopy category
	 of strong completion towers satisfying the same square-zero
	 and Mittag-Leffler hypotheses. In a curved CDG realization the
	 equivalence is the corresponding second-kind coderived
	 equivalence.
\end{enumerate}
\end{theorem}

\begin{proof}
\emph{Completed coalgebra.}
Each quotient map $p_N\colon \cA_{\le N+1} \twoheadrightarrow
\cA_{\le N}$ is strict by axiom~(4), inducing a strict
morphism
$q_N\colon \barB^{\mathrm{ch}}(\cA_{\le N+1}) \to
\barB^{\mathrm{ch}}(\cA_{\le N})$
compatible with differential, coproduct, coaugmentation, and
curvature. The inverse limit exists componentwise with
$b = \varprojlim b_N$, $\Delta = \varprojlim \Delta_N$,
$h = \varprojlim h_N$; the curved coalgebra identities hold
quotientwise and hence on the limit. Continuity is formal for
the inverse-limit topology. Pronilpotence: the reduced
coproduct on bar-length-$\le m$ elements is nilpotent on each
finite quotient, so the limit is topologically pronilpotent.

\emph{MC assembly.} Define
$\widehat\tau := \varprojlim_N \tau_N$. By
Lemma~\ref{lem:degree-cutoff}, modulo $F^{N+1}\cA$ only degrees
$\le N$ contribute to $\widehat\tau \star \widehat\tau$. The
MC equation holds modulo $F^{N+1}$ (where it reduces to the
finite-stage equation for $\tau_N$), hence on the completed
coalgebra. This is claim~(2).

\emph{Completed cobar.} Each
$\Omega^{\mathrm{ch}}(\barB^{\mathrm{ch}}(\cA_{\le N}))$ is
a curved dg chiral algebra; compatible transition maps give a
complete curved dg chiral algebra in the limit, with continuous
operations inherited from the finite stages. This is claim~(3).

\emph{Counit.} The quotient of
$\widehat\epsilon$ modulo $F^{N+1}$ is the finite-stage counit
$\epsilon_N$, which is a quasi-isomorphism by
Theorem~\ref{thm:finite-type-bar-cobar}. The finite-window
hypothesis gives the Mittag-Leffler condition for the source
and target cohomology systems in Lemma~\ref{lem:filtered-comparison}.
Applying that lemma to the compatible morphisms
$\{\epsilon_N\}$ yields claim~(4).

\emph{Dual comparison.} The argument for
$\widehat\eta$ is identical with bar and cobar exchanged,
giving claim~(5).

\emph{Equivalence.} Claims~(4) and~(5)
are the two triangle identities up to quasi-isomorphism, so
completed bar and completed cobar form an adjoint equivalence
on the homotopy categories restricted to strong completion
towers. This is claim~(6).
\end{proof}

% ================================================================
% 7. MC4+ vs MC4-zero
% ================================================================

\section{Positive and Resonant Completion}
\label{sec:mc4}

The strong-completion theorem proves the formal MC4 mechanism once
the degree cutoff, square-zero or coderived ambient, and
Mittag-Leffler hypotheses are in place. The positive regime is the
case $\rho(\cA)=0$: every completed operation either strictly
raises the chosen positive filtration, or an auxiliary positive
grading gives the same finite-window statement.

\begin{corollary}[Positive MC4]
\label{cor:mc4-plus}
Let $\cA$ be a positive strong completion tower with
$\rho(\cA)=0$, or a tower equipped with a genuine auxiliary
positive grading whose bar differential satisfies the finite-window
cutoff and the square-zero or coderived ambient hypotheses of
Theorem~\ref{thm:mc4-strong}. Then the completed bar-cobar counit
is a quasi-isomorphism in the square-zero total model, and a
coacyclic-cone comparison in the curved CDG ambient.
\end{corollary}

\begin{proof}
The strong case is Theorem~\ref{thm:mc4-strong}. In the auxiliary
positive grading case, positivity supplies the finite-window
hypothesis of Proposition~\ref{prop:ml}, so the same inverse-limit
comparison applies. Finite-resonance towers with
$0<\rho(\cA)<\infty$ are not included in this corollary.
\end{proof}

\begin{definition}[Finite resonance]
\label{def:finite-resonance}
A filtered chiral algebra has finite resonance if the part of the
bar differential preserving filtration degree is controlled by a
finite-dimensional complex $R_\cA$. Its resonance rank is
\[
\rho(\cA)=\dim H^*(R_\cA).
\]
The equality $\rho=0$ is the positive case. The case
$0<\rho<\infty$ is a finite-dimensional obstruction problem
requiring explicit homotopy data.
\end{definition}

\begin{theorem}[Generic finite-resonance MC4${}^0$]
\label{thm:n4-mc4-zero-generic-parameters}
Let $\cA$ be Virasoro at generic central charge or a principal
$\cW_N$ algebra at generic shifted level, equipped with its
positive-energy weight-completed finite-resonance model. Assume the
resonance complex $R_\cA$ is finite dimensional and is represented
inside the completed bar complex by either a chain-level strong
deformation retract or a finite-window stabilization. Then the
MC4${}^0$ completed bar-cobar comparison holds in the same
square-zero, respectively coderived, ambient as
Theorem~\ref{thm:mc4-strong}.
\end{theorem}

\begin{proof}
The positive-energy filtration decomposes the completed bar complex
into finite weight windows. The part of the differential preserving
filtration is, by hypothesis, represented by the finite complex
$R_\cA$; the remaining summands strictly raise weight and therefore
satisfy the cutoff of Proposition~\ref{prop:ml}. The chosen strong
deformation retract, or equivalently the finite-window
stabilization, identifies the finite-resonance summand with a fixed
finite complex at every sufficiently large stage. Hence the inverse
system of finite-stage cohomologies is Mittag-Leffler: the positive
part stabilizes weight by weight, and the resonance part is finite
dimensional. The completed comparison is then the inverse-limit
comparison of Theorem~\ref{thm:mc4-strong}. The coderived statement
is obtained by replacing the square-zero quasi-isomorphism with the
same finite-stage coacyclic-cone comparison before taking the
Mittag-Leffler limit.
\end{proof}

\begin{proposition}[Screening obstruction for the principal
$\cW_N$ route]
\label{prop:ff-screening-coproduct-obstruction}
Let $\mathfrak h$ be the rank-$(N-1)$ Heisenberg vertex algebra
with primitive chiral coproduct $\Delta_z^{\mathfrak h}$, and let
$Q_{\alpha_i}$ be the Feigin--Frenkel screening charge
\cite{FeiginFrenkel92} at shifted
level~$\Psi$. A strong deformation retract presenting
\[
\cW_N=\bigcap_i\ker Q_{\alpha_i}\subset \mathfrak h
\]
must account for the chiral $1$-cocycle
\[
[Q_{\alpha_i},\Delta_z^{\mathfrak h}]
=
\frac{\Psi-1}{\Psi}\,
\bigl(:e^{\alpha_i\cdot\varphi(0)}:\bigr)\otimes
\bigl(:e^{\alpha_i\cdot\varphi(-z)}:\bigr)\cdot(1+O(z)).
\]
At $\Psi=1$ this obstruction vanishes. At generic $\Psi$ it is
nonzero and must vanish under the homotopy data or be absorbed in a
completed coderived ambient.
\end{proposition}

\begin{proof}
This is the Vol.~I screening-coproduct computation: expanding the
primitive Heisenberg coproduct on the mode decomposition of
$:e^{\alpha_i\cdot\varphi(\zeta)}:$ gives the one-loop contraction
$\alpha_i^2\log z=(2/\Psi)\log z$, whose regularized residue is
$(\Psi-1)/\Psi$. The same coefficient is the Miura cross-term
coefficient for the spin-$s$ principal $\cW_N$ generators.
\end{proof}

\section{Standard Families}
\label{sec:examples}

\subsection{$\beta\gamma$ and affine Kac-Moody}

For the $\beta\gamma_\lambda$ system the conformal weights are
$\lambda$ and $1-\lambda$; the physical $\beta\gamma$ point
$\lambda=1$ gives weights $1$ and $0$. The raw conformal filtration
therefore is not a positive strong completion tower. Its shadow
data are contact class~$C$:
\[
S_3=0,\qquad S_4=-\frac{5}{12},\qquad S_r=0\quad(r\ge 5).
\]
Any MC4 application must use an auxiliary filtration that proves
the relevant bar-level cutoff.

For $V_k(\fg)$ at non-critical level the conformal filtration gives
the finite-window bar cutoff of Proposition~\ref{prop:examples}.
The algebra-level strong filtration fails because the
$J^a_{(0)}J^b$ bracket has weight~$1$ from two weight-$1$ inputs.
The critical
level $k=-h^\vee$ lies outside the finite-stage proved regime
because the Feigin--Frenkel centre changes the completion problem.

\subsection{Virasoro and principal $\cW$}

Virasoro has one weight-$2$ generator and OPE
\[
T(z)T(w)\sim\frac{c/2}{(z-w)^4}
+\frac{2T(w)}{(z-w)^2}
+\frac{\partial T(w)}{z-w}.
\]
The central fourth-order pole is the finite resonance direction.
Generic Virasoro MC4${}^0$ is therefore the finite-resonance
case of Theorem~\ref{thm:n4-mc4-zero-generic-parameters}.

For principal $\cW^k(\fg,f_{\mathrm{prin}})$ the generator weights
are the exponents $m_i+1$. In type~$A_{N-1}$ these are
$2,\ldots,N$. The standard principal-stage $W_\infty$ tower has a
proved bar-level weight-cutoff route when the quotient system and
square-zero/coderived hypotheses are specified. A generic
Feigin--Frenkel screening descent for individual principal
$\cW_N$ stages must still confront
Proposition~\ref{prop:ff-screening-coproduct-obstruction}.

\begin{center}
\begin{tabular}{lccc}
\toprule
Algebra & Input & Conclusion & Regime \\
\midrule
$V_k(\fg)$, $k\ne-h^\vee$ & bar-level cutoff & proved with hypotheses & MC4${}^+$ \\
$V_\Lambda$ & bar-level cutoff & proved with hypotheses & MC4${}^+$ \\
$\beta\gamma_\lambda$ & auxiliary contact cutoff & requires cutoff & class $C$ \\
$\Vir_c$ generic & finite resonance model & proved with hypotheses & MC4${}^0$ \\
principal $\cW_N$ generic & finite resonance model & proved with hypotheses & MC4${}^0$ \\
principal $\cW_N$ screening presentation & screening SDR & requires SDR & descent \\
$V_{-h^\vee}(\fg)$ & Feigin--Frenkel centre & excluded & critical \\
\bottomrule
\end{tabular}
\end{center}

% ================================================================
% References
% ================================================================

\begin{thebibliography}{99}

\bibitem{BD04}
A.~Beilinson and V.~Drinfeld,
\emph{Chiral algebras},
American Mathematical Society Colloquium Publications, vol.~51,
American Mathematical Society, Providence, RI, 2004.

\bibitem{BGS}
A.~Beilinson, V.~Ginzburg, and W.~Soergel,
\emph{Koszul duality patterns in representation theory},
J.\ Amer.\ Math.\ Soc.\ \textbf{9} (1996), 473--527.

\bibitem{FeiginFrenkel92}
B.~Feigin and E.~Frenkel,
\emph{Affine Kac--Moody algebras at the critical level and
Gelfand--Dikii algebras},
Internat.\ J.\ Modern Phys.\ A \textbf{7} (1992), suppl.~1A,
197--215.

\bibitem{LV12}
J.-L.~Loday and B.~Vallette,
\emph{Algebraic operads},
Grundlehren der mathematischen Wissenschaften, vol.~346,
Springer, Heidelberg, 2012.

\bibitem{LorgatI}
R.~Lorgat,
\emph{Modular Koszul Duality, Volume I},
Monograph, 2026.

\bibitem{Positselski11}
L.~Positselski,
\emph{Two kinds of derived categories, Koszul duality, and
comodule-contramodule correspondence},
Mem.\ Amer.\ Math.\ Soc.\ \textbf{212} (2011), no.~996.

\bibitem{Priddy70}
S.~B.~Priddy,
\emph{Koszul resolutions},
Trans.\ Amer.\ Math.\ Soc.\ \textbf{152} (1970), 39--60.

\end{thebibliography}

\end{document}
